\section{Banach空间}


\subsection{无限维Banach空间的基}
\begin{definition}
	设$X$是无限维的Banach空间，$X$中的点列$\{e_n\}$如果满足：$X$中的任一元素$x$可唯一地表示成
	\begin{equation*}
		x=\sum_{n=1}^{+\infty}\xi_ne_n
	\end{equation*}
	其中$\xi_n$为实数或复数，仅与$x$有关，且右端的级数依$X$中的范数收敛，则称$\{e_n\}$为$X$的Schauder基或简称基。
\end{definition}
\begin{theorem}
	具有Schauder基的Banach空间可分。
\end{theorem}
\begin{proof}
	设$X$为具有Schauder基$\{e_n\}$的Banach空间。取$X$的一个子集$E$，$E$是以下点构成的集合：
	\begin{equation*}
		\left\{\sum_{i=1}^nr_ie_i:n\in\mathbb{N}^+,r_i\in\mathbb{Q}\right\}
	\end{equation*}
	任取$x\in X$，则$x$可表示为$\sum\limits_{n=1}^{+\infty}\xi_ie_i$。在$E$中取点列$\{x_n\}$，其中:
	\begin{gather*}
		x_n=\sum_{i=1}^{n}\eta_{in}e_i \\
		\eta_{in}\in\mathbb{Q},\;\lim_{n\to+\infty}\eta_{in}=\xi_i
	\end{gather*}
	因为：
	\begin{align*}
		||x_n-x||
		&=\left\|\sum_{i=1}^{n}(\eta_{in}-\xi_i)e_i-\sum_{i=n+1}^{+\infty}\xi_ie_i\right\| \\
		&\leqslant\left\|\sum_{i=1}^{n}(\eta_{in}-\xi_i)e_i\right\|+\left\|\sum_{i=n+1}^{+\infty}\xi_ie_i\right\|
	\end{align*}
	因为$\sum\limits_{n=1}^{+\infty}\xi_ne_n$依$X$中的范数收敛，所以后一项由Cauchy收敛准则趋于$0$，而前一项显然趋于$0$，所以$\{x_n\}\rightarrow x$，且$x_n\in E,\;\forall\;n\in\mathbb{N}^+$。因此$x$是$E$的一个聚点。由$x$的任意性，$X\subset\overline{E}$。而$E$是可数集合，因此$X$可分。
\end{proof}

\begin{corollary}
	设内积空间$X$中存在完备的规范正交系，则$X$可分。
\end{corollary}
\begin{proof}
	设$\{e_n\}$是$X$中完备的规范正交系，由\cref{theo:OthonormalSystem}(2)可知对任意的$x\in X$，有：
	\begin{equation*}
		x=\sum_{n=1}^{+\infty}(x,e_n)e_n
	\end{equation*}
	根据\cref{theo:Density}(3)可知由$\{e_n\}$张成的子空间在$X$中稠密，由\info{可列个可列是可列}可知$\{e_n\}$以有理数为系数的所有可能的线性组合构成的集在$X$中也稠密且可列，于是$X$可分。
\end{proof}
\begin{theorem}
	可分的内积空间$X$必存在完备的规范正交系。
\end{theorem}
\begin{proof}
	因为$X$可分，所以$X$中有可列的稠密子集$\{x_n\}$。由Schmidt正交化定理，利用$\{x_n\}$可以建立规范正交系$\{e_n\}$，且$\{x_n\}$与$\{e_n\}$张成的子空间相同。因为$\{x_n\}$在$X$中稠密，所以任取$y\in X$，存在$\{x_n\}$中的点列$\{y_n\}$使得$\{y_n\}\to y$。因为$x_n$可由$\{e_n\}$线性表出，所以$\{y_n\}$可由$\{e_n\}$线性表出，于是$\{y_n\}\subset\operatorname{span}\{e_n,\;n\in\mathbb{N}^+\}$。把$y_n$用$\{e_n\}$线性表出，就可以得到：
	\begin{equation*}
		y_k=\sum_{n=1}^{+\infty}(y_k,e_n)e_n,\;k\in\mathbb{N}^+
	\end{equation*}
	因为：
	\begin{align*}
		\left\|y-\sum_{n=1}^{+\infty}(y,e_n)e_n\right\|
		&\leqslant||y-y_k||+\left\|y_k-\sum_{n=1}^{+\infty}(y,e_n)e_n\right\| \\
		&=||y-y_k||+\left\|\sum_{n=1}^{+\infty}(y_k,e_n)e_n-\sum_{n=1}^{+\infty}(y,e_n)e_n\right\| \\
		&=||y-y_k||+\left\|\sum_{n=1}^{+\infty}(y_k-y,e_n)e_n\right\|
	\end{align*}
	又：
	\begin{align*}
		\left\|\sum_{n=1}^{+\infty}(y_k-y,e_n)e_n\right\|
		&=\sqrt{\left(\sum_{n=1}^{+\infty}(y_k-y,e_n)e_n,\sum_{n=1}^{+\infty}(y_k-y,e_n)e_n\right)} \\
		&=\sqrt{\sum_{n=1}^{+\infty}(y_k-y,e_n)\overline{(y_k-y,e_n)}}  
		\\
		&=\sqrt{\sum_{n=1}^{+\infty}|(y_k-y,e_n)|^2}
	\end{align*}
	由Bessel inequality可得：
	\begin{equation*}
		\left\|\sum_{n=1}^{+\infty}(y_k-y,e_n)e_n\right\|
		=\sqrt{\sum_{n=1}^{+\infty}|(y_k-y,e_n)|^2}
		\leqslant||y_k-y||
	\end{equation*}
	于是有：
	\begin{equation*}
		\left\|y-\sum_{n=1}^{+\infty}(y,e_n)e_n\right\|
		\leqslant2||y-y_k||\to0\quad(k\to+\infty)
	\end{equation*}
	即：
	\begin{equation*}
		\left\|y-\sum_{n=1}^{+\infty}(y,e_n)e_n\right\|=0w
	\end{equation*}
	所以：
	\begin{equation*}
		y=\sum_{n=1}^{+\infty}(y,e_n)e_n
	\end{equation*}
	由$y$的任意性，$\{e_n\}$是完备的。
\end{proof}
\subsubsection{可分Hilbert空间的同构性}
\begin{theorem}
	实（或复）可分Hilbert空间与实（或复）$l^2$空间等距同构，所以所有实（或复）可分Hilbert空间彼此等距同构。
\end{theorem}
\begin{proof}
	设$X$为实（或复）可分Hilbert空间，$\{e_n\}$是$X$中的一个完备规范正交系。任取$x\in X$，$\{c_n\}$是$x$关于$\{e_n\}$的Fourier系数集。作$X$到$l^2$中的映射：
	\begin{equation*}
		T:Tx=\{c_n\}
	\end{equation*}\par
	(1)由Parseval's identity，$T$是双射。\par
	(2)任取$x,y\in X$，$\alpha,\beta$为$X$对应数域中的两个数，则：
	\begin{equation*}
		T(\alpha x+\beta y)=\{(\alpha x+\beta y,e_n)\}=\{\alpha(x,e_n)+\beta(y,e_n)\}=\alpha Tx+\beta Ty
	\end{equation*}\par
	综上两点，$T$是$X$到$l^2$中的同构映射。\par
	任取$x,y\in X$，由$\{e_n\}$的完备性可得：
	\begin{align*}
		||x-y||^2
		&=\left(\sum_{n=1}^{+\infty}(x-y,e_n)e_n,\sum_{n=1}^{+\infty}(x-y,e_n)e_n\right) \\
		&=\sum_{n=1}^{+\infty}(x-y,e_n)\overline{(x-y,e_n)} \\
		&=\sum_{n=1}^{+\infty}|(x-y,e_n)|^2 \\
		&=||Tx-Ty||^2
	\end{align*}
	所以$||x-y||=||Tx-Ty||$，即$X$与$l^2$空间等距。\par
	综上，$X$与$l^2$等距同构。由$X$的任意性以及等距同构的传递性，所有实（或复）可分Hilbert空间彼此等距同构。
\end{proof}